% answers_short.tex —— 2008 年·普通高等学校招生考试·福建卷·文科 · 简短答案速查
% 由 tools/make_short_answers.py 自 answers.tex 的 \gksolhead/\gkanswer 提取，
% 并与 2008/2008_数学真题_福建文_答案要点.md 逐题比对后生成，勿手改（改源 MD 后重跑）。

\begin{gknotice}
\textbf{使用说明}：本节只给最终答案便于快速核对，完整推导见解析版。答案为编写者独立推导，\textbf{非官方标准答案}。
\end{gknotice}

\gkpart{选择题}
\begin{center}\begin{tabular}{ccccccccccccc}
\toprule
题号 & 1 & 2 & 3 & 4 & 5 & 6 & 7 & 8 & 9 & 10 & 11 & 12 \\
\midrule
答案 & A & C & C & B & C & D & A & A & A & D & A & B \\
\bottomrule
\end{tabular}\end{center}

\gkpart{填空题}
\begin{itemize}[leftmargin=2.6em]
  \item[\textbf{13.}] 0。
  \item[\textbf{14.}] $(-\infty,0)\cup(10,+\infty)$。
  \item[\textbf{15.}] $9\pi$。
  \item[\textbf{16.}] ①④。
\end{itemize}

\gkpart{解答题}
\begin{itemize}[leftmargin=2.6em]
  \item[\textbf{17.}] （1）$\tan A=2$；（2）值域 $\left[-1-\dfrac{4\sqrt5}5,\dfrac75\right]$。
  \item[\textbf{18.}] （1）$\dfrac3{20}$；（2）“密码被破译”的概率大。
  \item[\textbf{19.}] （1）证明见解析；（2）$\dfrac{\sqrt6}3$；（3）$\dfrac{2\sqrt3}3$。
  \item[\textbf{20.}] （1）$a_n=n$；（2）证明见解析。
  \item[\textbf{21.}] （1）$m=-3$，$n=0$；增区间 $(-\infty,0)$、$(2,+\infty)$，减区间 $(0,2)$。（2）$0<a<1$ 时极大值为 $-2$；$1<a<3$ 时极小值为 $-6$；$a=1$ 或 $a\geqslant3$ 时无极值。
  \item[\textbf{22.}] （1）$\dfrac{x^2}4+\dfrac{y^2}3=1$；（2）①证明见解析；②$\dfrac92$。
\end{itemize}
